% ============================================================
\chapter{The Modelling Process}
\label{ch:process}
% ============================================================

Galileo wrote that the book of nature is written in the language of
mathematics. But between nature and equations, there is a bridge to build:
the \emph{model}. A model is a deliberate simplification of reality that
captures the essential mechanisms while ignoring superfluous details. The art
of modelling --- for it is as much an art as a science --- consists in
choosing the right level of abstraction: enough detail for the model to be
useful, enough simplicity for it to be tractable.

\begin{center}
\textit{``Science is built of facts the way a house is built of bricks;
but an accumulation of facts is no more a science than a pile of
bricks is a house.''}\\[4pt]
--- Henri Poincar\'e
\end{center}

\bigskip

This opening chapter introduces the general approach to mathematical
modelling.  We describe the key steps from observing a phenomenon to
validating the model, and classify the main families of mathematical models.

% ------------------------------------------------------------
\section{What Is a Mathematical Model?}
% ------------------------------------------------------------

\begin{definition}[Mathematical model]
A \textbf{mathematical model} is a formal representation of a real-world
system (physical, biological, economic, etc.)\ using mathematical
objects (equations, inequalities, graphs, probability distributions)
that capture the essential relationships between the system's quantities.
\end{definition}

\begin{remark}
A model is never unique: the same phenomenon can be described by
different models of varying complexity.  The choice depends on the
question asked and the data available.
\end{remark}

% ------------------------------------------------------------
\section{Steps in the Modelling Process}
\label{sec:steps}
% ------------------------------------------------------------

The modelling process consists of five broad stages, forming an
iterative cycle.

\begin{enumerate}
  \item \textbf{Observation and problem identification.}
    Delineate the phenomenon, identify relevant quantities (state
    variables, parameters, inputs/outputs), and gather available data.

  \item \textbf{Formulation of assumptions.}
    State explicitly which interactions are neglected, which quantities
    are assumed constant, and what the spatial and temporal domain of
    validity is.

  \item \textbf{Mathematical formulation.}
    Translate assumptions into equations.  This may yield:
    \begin{itemize}
      \item an ordinary differential equation (ODE),
      \item a partial differential equation (PDE),
      \item an algebraic system,
      \item an optimisation problem,
      \item a stochastic model.
    \end{itemize}

  \item \textbf{Solution and analysis.}
    Solve the model analytically (when possible) or numerically.
    Study qualitative behaviour: existence and uniqueness, stability
    of equilibria, parameter dependence.

  \item \textbf{Validation and interpretation.}
    Compare model predictions with experimental data.  If the agreement
    is insufficient, return to Step~2 to revise the assumptions.
\end{enumerate}

\begin{intuition}[title=\textbf{The modelling cycle}]
Modelling is not a linear process but a \textbf{cycle}.  After
validation, one often revisits the assumptions, enriches the model,
or reduces its complexity.  A good modeller goes through several
iterations.
\end{intuition}

% ------------------------------------------------------------
\section{Classification of Models}
\label{sec:classification}
% ------------------------------------------------------------

\begin{definition}[Types of models]
Models can be classified along several axes:
\begin{enumerate}
  \item \textbf{Deterministic vs.\ stochastic.}
    A deterministic model predicts a unique trajectory for given initial
    conditions; a stochastic model incorporates randomness.
  \item \textbf{Continuous vs.\ discrete.}
    Time and/or space may vary continuously or in discrete steps.
  \item \textbf{Static vs.\ dynamic.}
    Static models describe equilibrium states; dynamic models describe
    temporal evolution.
  \item \textbf{Linear vs.\ nonlinear.}
    Linearity of the relationships greatly affects the analysis methods.
  \item \textbf{Parametric vs.\ non-parametric.}
    Whether the model structure is fixed a priori or learned from data.
\end{enumerate}
\end{definition}

\begin{example}
Consider a falling object.  Neglecting air resistance, the model is:
\[
  m\ddot{y} = -mg, \qquad y(0)=h,\quad \dot{y}(0)=0.
\]
This is a deterministic, continuous, dynamic, linear model.
Its solution is $y(t)=h-\tfrac{1}{2}gt^2$.

Adding quadratic air resistance yields a \textbf{nonlinear} model:
\[
  m\ddot{y} = -mg + k\dot{y}^2.
\]
\end{example}

% ------------------------------------------------------------
\section{Fundamental Construction Principles}
\label{sec:principles}
% ------------------------------------------------------------

\subsection{Conservation Laws}

Many models rest on \textbf{conservation laws}: conservation of mass,
energy, momentum, or number of individuals.

\begin{theorem}[Generic balance equation]
For a quantity $Q(t)$ contained in a domain $\Omega$, the generic
conservation law reads:
\[
  \frac{dQ}{dt} = (\text{inflow}) - (\text{outflow})
                 + (\text{source}) - (\text{sink}).
\]
\end{theorem}

\begin{example}
Let $N(t)$ be the size of a population with per-capita birth rate $b$
and death rate $d$.  The balance gives:
\[
  \frac{dN}{dt} = bN - dN = (b-d)N.
\]
This is the Malthus model, studied in detail in
Chapter~\ref{ch:populations}.
\end{example}

\subsection{Constitutive Laws}

Constitutive laws relate fluxes to state variables.

\begin{example}
\begin{itemize}
  \item \textbf{Fourier's law:} heat flux is proportional to the
    temperature gradient: $\mathbf{q}=-\kappa\nabla T$.
  \item \textbf{Fick's law:} diffusive flux is proportional to the
    concentration gradient: $\mathbf{J}=-D\nabla c$.
  \item \textbf{Hooke's law:} stress is proportional to strain:
    $\sigma = E\varepsilon$.
\end{itemize}
\end{example}

\subsection{Variational Principles}

Some models are obtained by minimising (or making stationary) a functional.

\begin{definition}[Action functional]
In mechanics, the principle of least action states that the true
trajectory makes stationary the functional
\[
  \mathcal{S}[q] = \int_{t_1}^{t_2} L(q,\dot{q},t)\,dt,
\]
where $L$ is the Lagrangian.  The resulting Euler--Lagrange equations
will be studied in Chapter~\ref{ch:optimization}.
\end{definition}

% ------------------------------------------------------------
\section{Qualities of a Good Model}
% ------------------------------------------------------------

\begin{bestpractice}[title=\textbf{Quality criteria}]
A good model strikes a balance between:
\begin{itemize}
  \item \textbf{Simplicity} (parsimony): the number of parameters should
    remain reasonable (Occam's razor).
  \item \textbf{Accuracy}: predictions must be close enough to
    observations.
  \item \textbf{Robustness}: small parameter changes should not cause
    dramatic prediction shifts.
  \item \textbf{Interpretability}: parameters should have a physical
    or biological meaning.
\end{itemize}
\end{bestpractice}

\begin{proposition}[Parsimony principle]
Given two models of comparable accuracy, one prefers the one with
fewer free parameters.  This is formalised by information criteria
such as the AIC (Akaike Information Criterion):
\[
  \mathrm{AIC} = 2k - 2\ln \hat{L},
\]
where $k$ is the number of parameters and $\hat{L}$ the maximum
likelihood.
\end{proposition}

% ------------------------------------------------------------
\section{Sensitivity and Uncertainty}
\label{sec:sensitivity}
% ------------------------------------------------------------

\begin{definition}[Sensitivity analysis]
Sensitivity analysis studies how the output $y$ of a model varies
when a parameter $p$ changes.  The \textbf{local sensitivity index}
is defined by:
\[
  S_p = \frac{p}{y}\,\frac{\partial y}{\partial p}.
\]
An index $\abs{S_p}\gg 1$ signals high sensitivity to $p$.
\end{definition}

\begin{example}
For the exponential model $y(t)=y_0 e^{rt}$:
\[
  S_r = \frac{r}{y}\,\frac{\partial y}{\partial r}
      = \frac{r}{y_0 e^{rt}}\cdot y_0 t e^{rt} = rt.
\]
Sensitivity to $r$ grows linearly with time: small errors in $r$
lead to large errors in $y$ over long horizons.
\end{example}

\begin{codeblock}[title=\textbf{Sensitivity analysis in Python}]
\begin{minted}{python}
import numpy as np
import matplotlib.pyplot as plt

def model(t, r, y0=1.0):
    return y0 * np.exp(r * t)

t = np.linspace(0, 10, 200)
r_nom = 0.5
dr = 0.05  # 10% perturbation

y_nom = model(t, r_nom)
y_plus = model(t, r_nom + dr)
y_minus = model(t, r_nom - dr)

fig, (ax1, ax2) = plt.subplots(1, 2, figsize=(12, 4))
ax1.plot(t, y_nom, 'b-', label=f'$r = {r_nom}$')
ax1.fill_between(t, y_minus, y_plus, alpha=0.3,
                 label=f'$r = {r_nom} \\pm {dr}$')
ax1.set_xlabel('$t$'); ax1.set_ylabel('$y(t)$')
ax1.legend(); ax1.set_title('Solutions')

S_r = r_nom * t
ax2.plot(t, S_r, 'r-')
ax2.set_xlabel('$t$'); ax2.set_ylabel('$S_r$')
ax2.set_title('Sensitivity index $S_r = rt$')
plt.tight_layout(); plt.savefig('sensitivity.pdf')
\end{minted}
\end{codeblock}

% ------------------------------------------------------------
\section{Fitting Models to Data}
\label{sec:fitting}
% ------------------------------------------------------------

Once a model is built, its parameters must be estimated from data.

\begin{definition}[Least squares]
Given observations $(t_i, y_i^{\mathrm{obs}})$, $i=1,\ldots,n$, and a
model $y(t;\mathbf{p})$ depending on a parameter vector $\mathbf{p}$,
the \textbf{least-squares estimator} is:
\[
  \hat{\mathbf{p}} = \arg\min_{\mathbf{p}}
    \sum_{i=1}^{n} \bigl(y_i^{\mathrm{obs}} - y(t_i;\mathbf{p})\bigr)^2.
\]
\end{definition}

\begin{codeblock}[title=\textbf{Least-squares fitting}]
\begin{minted}{python}
import numpy as np
from scipy.optimize import curve_fit

np.random.seed(42)
t_data = np.linspace(0, 5, 20)
r_true, y0_true = 0.8, 2.0
y_data = y0_true * np.exp(r_true * t_data) \
         + np.random.normal(0, 1, len(t_data))

def model_func(t, y0, r):
    return y0 * np.exp(r * t)

popt, pcov = curve_fit(model_func, t_data, y_data, p0=[1, 0.5])
print(f"Estimated: y0 = {popt[0]:.3f}, r = {popt[1]:.3f}")
print(f"Std errors: {np.sqrt(np.diag(pcov))}")
\end{minted}
\end{codeblock}

% ------------------------------------------------------------
\section{Complete Example: Cooling of a Cup of Coffee}
\label{sec:coffee}
% ------------------------------------------------------------

Let us illustrate the full process on a classic example.

\paragraph{Step 1: Observation.}
A hot cup of coffee placed on a table cools down over time.  We measure
the temperature $T(t)$.

\paragraph{Step 2: Assumptions.}
\begin{itemize}
  \item The coffee temperature is spatially uniform (lumped model).
  \item Ambient temperature $T_a$ is constant.
  \item Heat flux is proportional to the temperature difference
    (Newton's law of cooling).
\end{itemize}

\paragraph{Step 3: Model.}
Newton's law gives:
\[
  \frac{dT}{dt} = -k(T - T_a), \qquad T(0)=T_0,
\]
where $k>0$ is the heat transfer coefficient.

\paragraph{Step 4: Solution.}
This is a first-order linear ODE with solution:
\[
  T(t) = T_a + (T_0 - T_a)\,e^{-kt}.
\]

\begin{keyformulas}[title=\textbf{Newton's Law of Cooling}]
\[
  T(t) = T_a + (T_0 - T_a)\,e^{-kt}, \qquad k > 0.
\]
\begin{itemize}
  \item $T_0$: initial temperature.
  \item $T_a$: ambient temperature.
  \item $k$: cooling constant ($\mathrm{s}^{-1}$).
  \item Half-cooling time: $t_{1/2} = \ln 2 / k$.
\end{itemize}
\end{keyformulas}

\paragraph{Step 5: Validation.}

\begin{codeblock}[title=\textbf{Simulation and fit --- cooling}]
\begin{minted}{python}
import numpy as np
import matplotlib.pyplot as plt
from scipy.optimize import curve_fit

np.random.seed(0)
t_exp = np.arange(0, 31, 2)
T_a, T0_true, k_true = 22.0, 90.0, 0.07
T_exp = T_a + (T0_true - T_a) * np.exp(-k_true * t_exp) \
        + np.random.normal(0, 1.5, len(t_exp))

def newton_cooling(t, T0, k):
    return T_a + (T0 - T_a) * np.exp(-k * t)

popt, pcov = curve_fit(newton_cooling, t_exp, T_exp, p0=[85, 0.05])
print(f"T0 = {popt[0]:.1f} C,  k = {popt[1]:.4f} /min")

t_fine = np.linspace(0, 30, 300)
plt.figure(figsize=(8, 4))
plt.scatter(t_exp, T_exp, c='red', label='Data')
plt.plot(t_fine, newton_cooling(t_fine, *popt), 'b-',
         label='Fitted model')
plt.axhline(T_a, color='gray', ls='--', label=f'$T_a = {T_a}$ C')
plt.xlabel('Time (min)'); plt.ylabel('Temperature (C)')
plt.legend(); plt.title("Newton's Cooling Law")
plt.tight_layout(); plt.savefig('cooling.pdf')
\end{minted}
\end{codeblock}

% ------------------------------------------------------------
\section{Limitations and Common Pitfalls}
% ------------------------------------------------------------

\begin{warning}[title=\textbf{Pitfalls to avoid}]
\begin{enumerate}
  \item \textbf{Overfitting.}
    A model with too many parameters may fit the training data
    perfectly but fail at prediction.
  \item \textbf{Reckless extrapolation.}
    A model validated on $[0,T]$ should not be extrapolated to
    $t\gg T$ without caution.
  \item \textbf{Correlation vs.\ causation.}
    A good fit does not prove a causal link.
  \item \textbf{Forgetting units.}
    Never add quantities with different dimensions
    (see Chapter~\ref{ch:dimensional}).
\end{enumerate}
\end{warning}

% ------------------------------------------------------------
\section{Exercises}
% ------------------------------------------------------------

\begin{exercise}
A tank initially contains $V_0=1000$~L of pure water.  A salt solution
of concentration $c_{\mathrm{in}}=30$~g/L flows in at rate
$q_{\mathrm{in}}=5$~L/min.  The well-mixed solution flows out at rate
$q_{\mathrm{out}}=5$~L/min.
\begin{enumerate}
  \item Write the ODE for the amount of salt $Q(t)$ in the tank.
  \item Solve the ODE and find the steady-state concentration.
  \item Plot $Q(t)$ for $t\in[0,400]$~min using Python.
  \item Discuss what changes if $q_{\mathrm{out}}\neq q_{\mathrm{in}}$.
\end{enumerate}
\end{exercise}

\begin{exercise}
Consider the model $y(t)=A\sin(\omega t+\varphi)$ for a periodic signal.
\begin{enumerate}
  \item Generate synthetic noisy data with $A=3$, $\omega=2\pi/5$,
    $\varphi=\pi/4$, and Gaussian noise of standard deviation $0.5$.
  \item Fit the model by least squares.
  \item Compute 95\% confidence intervals for $A$, $\omega$, $\varphi$.
\end{enumerate}
\end{exercise}

\begin{exercise}
For Newton's cooling model:
\begin{enumerate}
  \item Show that $T(t)$ is strictly decreasing if $T_0>T_a$ and
    strictly increasing if $T_0<T_a$.
  \item Compute the sensitivity index $S_k$ and interpret.
  \item Modify the model to account for radiation (Stefan--Boltzmann:
    flux $\propto T^4-T_a^4$).  Solve numerically and compare.
\end{enumerate}
\end{exercise}

\begin{exercise}
The number of atoms $N(t)$ in a radioactive substance satisfies
$dN/dt = -\lambda N$.
\begin{enumerate}
  \item Show that the half-life is $t_{1/2}=\ln 2/\lambda$.
  \item Write a Python program that estimates $\lambda$ from data
    $(t_i,N_i)$ by (a) nonlinear least squares and (b) linear
    regression on $\ln N$.
  \item Compare the two methods and discuss their merits.
\end{enumerate}
\end{exercise}
